% Part: first-order-logic
% Chapter: beyond
% Section: higher-order-logic

\documentclass[../../../include/open-logic-section]{subfiles}

\begin{document}

\olfileid{fol}{byd}{hol}

\olsection{উচ্চতর-ক্রমের যুক্তিবিদ্যা}

প্রথম-ক্রমের যুক্তিবিদ্যা থেকে দ্বিতীয়-ক্রমের যুক্তিবিদ্যায় যাওয়ায়
আনুষ্ঠানিক ভাষার মধ্যেই প্রথম-ক্রমীয় !!{domain}-এর বস্তুগুলির সেট নিয়ে
কথা বলা সম্ভব হয়েছে। সেখানেই থামব কেন? যেমন, তৃতীয়-ক্রমের
যুক্তিবিদ্যায় বস্তুর সেটগুলির সেট, অথবা বস্তু ও বস্তুর সেট উভয়ই ধারণ
করে এমন সেট নিয়েও কাজ করা সম্ভব হওয়া উচিত। আর চতুর্থ-ক্রমের
যুক্তিবিদ্যায় সেই ধরনের বস্তুর সেট নিয়ে কথা বলা যাবে। অনুমান করা যায়,
এই ধারণাটিকে ইচ্ছামতো পুনরাবৃত্ত করা যায়।

ব্যবহারিক ক্ষেত্রে উচ্চতর-ক্রমের যুক্তিবিদ্যা প্রায়ই সম্পর্কের বদলে
অপেক্ষকের সাহায্যে !!{formula}ted করা হয়। (স্বাভাবিক সনাক্তকরণগুলি
ধরলে এই পার্থক্যটি অপরিহার্য নয়।) কয়েকটি মৌলিক “জাতি” $A$, $B$,
$C$,~\dots নিয়ে---এখন এগুলিকে আমরা “টাইপ” বলব---এই নিয়মে নতুন টাইপ
তৈরি করা যায়:
\begin{quote}
$\sigma$ ও $\tau$ সসীম টাইপ হলে $\sigma \to \tau$-ও সসীম টাইপ।
\end{quote}
টাইপগুলিকে এমন বাক্যতাত্ত্বিক “চিহ্নপত্র” হিসেবে ভাবো যা আমাদের
!!{domain}-এ কাঙ্ক্ষিত বস্তুগুলিকে শ্রেণীবদ্ধ করে; $\sigma \to \tau$
দিয়ে এমন বস্তু বোঝায় যেগুলি $\sigma$ টাইপের বস্তু থেকে $\tau$ টাইপের
বস্তুতে অপেক্ষক। যেমন, “সত্য” ও “মিথ্যা” সত্যমানগুলির একটি টাইপ
$\Omega$ এবং স্বাভাবিক সংখ্যার একটি টাইপ $\Nat$ রাখতে চাইতে পারি।
সেক্ষেত্রে $\Nat \to \Omega$ টাইপের বস্তুগুলিকে এক-স্থানীয় সম্পর্ক
অথবা $\Nat$-এর উপসেট; $\Nat \to \Nat$ টাইপের বস্তুগুলিকে স্বাভাবিক
সংখ্যা থেকে স্বাভাবিক সংখ্যায় অপেক্ষক; এবং $(\Nat \to \Nat) \to
\Nat$ টাইপের বস্তুগুলিকে “ফাংশনাল” ভাবা যায়---অর্থাৎ, এমন
উচ্চতর-টাইপের অপেক্ষক যা অপেক্ষক নিয়ে সংখ্যা দেয়।

দ্বিতীয়-ক্রমের যুক্তিবিদ্যার মতো উচ্চতর-ক্রমের যুক্তিবিদ্যাকেও এক
ধরনের বহুজাতীয় যুক্তিবিদ্যা হিসেবে ভাবা যায়, যেখানে বিবেচ্য প্রত্যেক
টাইপের বস্তুর জন্য আলাদা একটি জাতি আছে। তবে উচ্চতর-টাইপের যুক্তিবিদ্যার
বাক্যতত্ত্বকে গোড়া থেকে সংজ্ঞায়িত করাই সাধারণত বেশি পরিষ্কার। যেমন,
সসীম টাইপের একটি সেটকে নিচের মতো আরোহক্রমে সংজ্ঞায়িত করা যায়:
\begin{enumerate}
\item $\Nat$ একটি সসীম টাইপ।
\item $\sigma$ ও $\tau$ সসীম টাইপ হলে $\sigma \to \tau$-ও সসীম
  টাইপ।
\item $\sigma$ ও $\tau$ সসীম টাইপ হলে $\sigma \times \tau$-ও সসীম
  টাইপ।
\end{enumerate}
স্বজ্ঞাতভাবে $\Nat$ স্বাভাবিক সংখ্যার টাইপ, $\sigma \to \tau$
$\sigma$ থেকে $\tau$-তে অপেক্ষকের টাইপ, এবং $\sigma \times \tau$
একটি $\sigma$ থেকে ও একটি $\tau$ থেকে নেওয়া বস্তুজোড়ার টাইপ নির্দেশ
করে। এরপর পদের একটি সেটকে নিচের মতো আরোহক্রমে সংজ্ঞায়িত করা যায়:
\begin{enumerate}
\item প্রত্যেক টাইপ $\sigma$-র জন্য $x$, $y$, $z$, \dots চলরাশির
  একটি ভাণ্ডার আছে; চলরাশিগুলি $\sigma$ টাইপের।
\item $\Obj 0$ হলো $\Nat$ টাইপের একটি পদ।
\item $\Obj S$ (উত্তরসূরি) হলো $\Nat \to \Nat$ টাইপের একটি পদ।
\item $s$ যদি $\sigma$ টাইপের পদ এবং $t$ যদি $\Nat \to (\sigma
  \to \sigma)$ টাইপের পদ হয়, তবে $\Obj{R}_{st}$ হলো $\Nat \to
  \sigma$ টাইপের একটি পদ।
\item $s$ যদি $\tau \to \sigma$ টাইপের পদ এবং $t$ যদি $\tau$ টাইপের
  পদ হয়, তবে $s(t)$ হলো $\sigma$ টাইপের একটি পদ।
\item $s$ যদি $\sigma$ টাইপের পদ এবং $x$ যদি $\tau$ টাইপের চলরাশি
  হয়, তবে $\lambd[x][s]$ হলো $\tau \to \sigma$ টাইপের একটি পদ।
\item $s$ যদি $\sigma$ টাইপের পদ এবং $t$ যদি $\tau$ টাইপের পদ হয়,
  তবে $\tuple{s, t}$ হলো $\sigma \times \tau$ টাইপের একটি পদ।
\item $s$ যদি $\sigma \times \tau$ টাইপের পদ হয়, তবে $p_1(s)$
  হলো $\sigma$ টাইপের এবং $p_2(s)$ হলো $\tau$ টাইপের পদ।
\end{enumerate}
স্বজ্ঞাতভাবে $\Obj{R}_{st}$ নিচের পুনরাবৃত্ত সংজ্ঞায় নির্ধারিত
অপেক্ষকটিকে নির্দেশ করে:
\begin{align*}
\Obj{R}_{st}(0) & = s \\
\Obj{R}_{st}(x+1) & = t(x, R_{st}(x)),
\end{align*}
$\tuple{s, t}$ সেই জোড়াটি নির্দেশ করে যার প্রথম উপাংশ~$s$ এবং দ্বিতীয়
উপাংশ~$t$; আর $p_1(s)$ ও~$p_2(s)$ যথাক্রমে~$s$-এর প্রথম ও দ্বিতীয়
উপাদান (“অভিক্ষেপ”) নির্দেশ করে। সব শেষে $\lambd[x][s]$ নিচের
সমীকরণে সংজ্ঞায়িত অপেক্ষক~$f$-কে নির্দেশ করে:
\[
f(x) = s
\]
$\tau$ টাইপের যেকোনো~$x$-এর জন্য। তাই দফা (6) পদ ব্যবহার করে অপেক্ষক
সংজ্ঞায়িত করার সুযোগ দিয়ে ধর্মনির্দেশের একটি রূপ দেয়।
% BN-SRC-111 TARGET-CORRECTION-BEGIN
\par\smallskip\noindent\textnormal{\small\textbf{উৎস-সংশোধন (BN-SRC-111):} হিমায়িত ইংরেজি উৎসে $\lambd[x][s]$-এর ব্যাখ্যার এই বাক্যে $x$-কে ভুল করে $\sigma$ টাইপের বলা হয়েছে, যদিও দফা (6)-এ $x$-এর টাইপ $\tau$ এবং ফল-পদ $s$-এর টাইপ $\sigma$। বাংলা পাঠে $x$-এর টাইপ $\tau$ রাখা হয়েছে।} % BN-SRC-111 TARGET-CORRECTION-END
একই টাইপের পদগুলির মধ্যকার !!{identity} বিবৃতি $\eq[s][t]$, পরিচিত
প্রস্তাবনামূলক সংযোজক এবং উচ্চতর-টাইপের পরিমাণায়ন থেকে
!!^{formula}s গড়ে ওঠে। এরপর উপরে সংজ্ঞায়িত পদগুলি পরিচালনাকারী মৌলিক
সমীকরণ এবং পরিমাণসূচক ও !!{identity}-সহ যুক্তিবিদ্যার পরিচিত বিধিগুলিকে
ব্যবস্থাটির স্বতঃসিদ্ধ হিসেবে নেওয়া যায়।

সসীম টাইপের ব্যবস্থায় সত্যমানের একটি টাইপ $\Omega$ যোগ করলে তার
ব্যবহার পরিচালনার জন্যও স্বতঃসিদ্ধ রাখতে হয়। বস্তুত যথেষ্ট কৌশলী হলে
জটিল !!{formula}s সম্পূর্ণ বাদ দিয়ে সেগুলির বদলে $\Omega$ টাইপের পদ
বসানো যায়!{} তখন প্রমাণ-ব্যবস্থাটিকেও সেইমতো পরিবর্তন করা যায়। এর
ফল মূলত ১৯৩০-এর দশকে অ্যালোনজো চার্চ প্রবর্তিত \emph{সরল
টাইপতত্ত্ব}।

দ্বিতীয়-ক্রমের যুক্তিবিদ্যার মতো উচ্চতর-টাইপের অর্থতত্ত্বেরও বিভিন্ন
সংস্করণ ব্যবহার করা যেতে পারে। পূর্ণ সংস্করণে $\sigma \to \tau$
টাইপের চলরাশিগুলি $\sigma$ টাইপের বস্তু থেকে $\tau$ টাইপের বস্তুতে
\emph{সমস্ত} অপেক্ষকের সেটের উপর বিস্তৃত হয়। প্রত্যাশামতোই এই
অর্থতত্ত্ব এত শক্তিশালী যে এর জন্য কোনো সম্পূর্ণ ও কার্যকর
!!{derivation} ব্যবস্থা থাকতে পারে না। তবে একটি দুর্বলতর অর্থতত্ত্ব
বিবেচনা করা যায়, যেখানে প্রত্যেক টাইপ~$\tau$-এর জন্য উপাদানের
সেট~$T_\tau$ এবং প্রয়োগ, অভিক্ষেপ ইত্যাদির উপযুক্ত ক্রিয়া নিয়ে
!!a{structure} গঠিত হয়। খুঁটিনাটি সঠিকভাবে সম্পন্ন করলে উপরে বর্ণিত
ধরনের !!{derivation} ব্যবস্থাগুলির জন্য সম্পূর্ণতা উপপাদ্য পাওয়া যায়।

উচ্চতর-টাইপের যুক্তিবিদ্যা আকর্ষণীয়, কারণ এটি এমন একটি কাঠামো দেয়
যার মধ্যে স্বাভাবিক উপায়ে গণিতের অনেকখানি নিবেশিত করা যায়: $\Nat$
থেকে শুরু করে বাস্তব সংখ্যা, অবিচ্ছিন্ন অপেক্ষক ইত্যাদি সংজ্ঞায়িত করা
যায়। স্বজ্ঞাবাদী যুক্তিবিদ্যার প্রসঙ্গেও এটি বিশেষ আকর্ষণীয়, কারণ
টাইপগুলির স্পষ্ট “নির্মাণমূলক” ব্যাখ্যা আছে। বস্তুত উচ্চতর-টাইপের
অর্থতত্ত্বের নির্মাণমূলক সংস্করণ (ধ্রুপদি যুক্তিবিদ্যার বদলে স্বজ্ঞাবাদী
যুক্তিবিদ্যার ভিত্তিতে) গড়ে তোলা যায়; সেগুলি এই নির্মাণমূলক
ব্যাখ্যাগুলিকে বেশ সুন্দরভাবে স্পষ্ট করে এবং অনেক দিক দিয়েই ধ্রুপদি
প্রতিরূপগুলির চেয়ে বেশি আকর্ষণীয়।

\end{document}
